\section{Gaussian quadratics over F{\o}lner sets}
\label{sec:gaussian}

We retain the algebra \(C\), its invariant commuting orbit algebras
\(C_i\), and the variables \(x_{i,g}\), \(z_{N,g}\), and \(k_{N;g,h}\)
from Section~\ref{sec:rows}. In particular, \(\T(C)\) is nonempty,
the \(x_{i,g}\) are uniformly bounded by \(M\), and
Proposition~\ref{prop:joint-clt} applies to arbitrary sequences of
extreme traces. All traces of finite-dimensional matrices in this
section are unnormalized. For a matrix \(B\), we write
\(\|B\|_{\HS}=(\Tr(B^*B))^{1/2}\) and
\(\|B\|_1=\Tr|B|\).

\begin{proposition}[Uniform approximately invariant unitaries]\label{prop:uniform-unitaries}
For every finite set \(H\subset G\), every \(\eps>0\), and every
integer \(m\geq1\), there are \(t>0\) and a self-adjoint element
\(Q\in C\) such that \(V=\exp(\iu tQ)\) satisfies
\begin{align}
 \sup_{\varphi\in\T(C)}
 \|\bar\alpha_h(V)-V\|_{2,\varphi}
 &<\eps &&(h\in H), \label{eq:uniform-unitary-invariance}\\
 \sup_{\varphi\in\T(C)}|\varphi(V^\ell)|
 &<\eps &&(1\leq\ell\leq m).
 \label{eq:uniform-unitary-moments}
\end{align}
\end{proposition}

We first record the matrix and probability estimates used in the
proof.

\begin{lemma}[Square-root perturbation bound]\label{lem:square-root-trace}
If \(S,T\) are positive matrices of the same size, then
\[
 \|S^{1/2}-T^{1/2}\|_{\HS}^{\,2}\leq\|S-T\|_1.
\]
\end{lemma}

\begin{proof}
Set \(a=S^{1/2}\), \(b=T^{1/2}\), \(d=a-b\), and
\(u=\operatorname{sign}(d)\), with \(\operatorname{sign}(0)=0\).
Since \(a^2-b^2=ad+db\), cyclicity of the trace gives
\[
 \Tr\bigl(u(a^2-b^2)\bigr)
 =\Tr\bigl(|d|(a+b)\bigr).
\]
In a basis diagonalizing \(d\), positivity gives
\(a_{jj}+b_{jj}\geq|a_{jj}-b_{jj}|=|d_j|\). Thus
\[
 \Tr\bigl(|d|(a+b)\bigr)
 =\sum_j |d_j|(a_{jj}+b_{jj})
 \geq\sum_j d_j^2=\Tr(d^2).
\]
On the other hand, \(u\) is a contraction, so trace-norm duality
gives
\[
 \Tr\bigl(u(a^2-b^2)\bigr)\leq\Tr|a^2-b^2|.
\]
These two inequalities prove the assertion.
\end{proof}

\begin{lemma}[Gaussian quadratic estimates]\label{lem:gaussian-quadratic}
Let \(W\) be a standard Gaussian vector in a finite-dimensional
real Euclidean space, and let \(B\) be a real symmetric operator
on that space. Put
\[
 q_B=\langle W,BW\rangle-\Tr(B).
\]
Then
\begin{equation}\label{eq:quadratic-variance-characteristic}
 \|q_B\|_{L^2}=\sqrt{2}\,\|B\|_{\HS},
 \qquad
 \bigl|\mathbb E\exp(\iu s q_B)\bigr|
 \leq\bigl(1+4s^2\|B\|_{\HS}^2\bigr)^{-1/4}
 \quad(s\in\R).
\end{equation}
Its law has an exponential moment of its absolute value and is
determined by its moments.
\end{lemma}

\begin{proof}
Orthogonal diagonalization of \(B\) gives
\[
 q_B\ \overset{\mathrm{law}}{=}
 \sum_{j=1}^d\lambda_j(U_j^2-1),
\]
where \(\lambda_j\) are its eigenvalues and \(U_j\) are independent
standard real normal variables. The summands are centered and
have variances \(2\lambda_j^2\), proving the first formula.
The one-dimensional Gaussian integral gives
\[
 \bigl|\mathbb E\exp(\iu s q_B)\bigr|
 =\prod_{j=1}^d(1+4s^2\lambda_j^2)^{-1/4}.
\]
For nonnegative numbers \(a_j\), one has
\(\prod_j(1+a_j)\geq1+\sum_j a_j\). Applying this inequality with
\(a_j=4s^2\lambda_j^2\) proves the second formula in
Equation~\eqref{eq:quadratic-variance-characteristic}.

Choose \(c>0\) so that \(2c\max_j|\lambda_j|<1\); if \(B=0\),
any \(c>0\) will do. The pointwise estimate
\[
 \exp\!\left(c\left|\sum_j\lambda_j(U_j^2-1)\right|\right)
 \leq
 \exp\!\left(c\sum_j|\lambda_j|\right)
 \prod_j\exp(c|\lambda_j|U_j^2)
\]
and the Gaussian integral show that
\(\mathbb E\exp(c|q_B|)<\infty\).

For completeness, let \(\nu\) be another probability measure with
the same moments as \(q_B\). By the nonnegative power series for
\(\cosh\) and monotone convergence,
\[
 \int_{\R}\cosh(cx)\,d\nu(x)
 =\sum_{r=0}^{\infty}
   \frac{c^{2r}}{(2r)!}\mathbb E(q_B^{2r})
 =\mathbb E\cosh(cq_B)<\infty.
\]
Since \(e^{c|x|}\leq2\cosh(cx)\), the measure \(\nu\) also has an
exponential moment. The characteristic functions of \(\nu\) and
of \(q_B\) extend analytically to the strip
\(\{z\in\C:|\operatorname{Im}z|<c\}\). Their Taylor coefficients
at zero agree, so the Identity Theorem makes these analytic
functions equal. Uniqueness of characteristic functions proves
that \(\nu\) is the law of \(q_B\).
\end{proof}

The purpose of square-root normalization is already visible for a
positive real symmetric matrix \(P\) with \(\Tr(P)=1\).
Equation~\eqref{eq:quadratic-variance-characteristic} gives
\[
 \operatorname{Var}(q_{P^{1/2}})=2,
 \qquad
 \operatorname{Var}(q_P)=2\Tr(P^2).
\]
The latter variance is \(2\) when \(P\) is a rank-one projection,
but \(2/n\) when \(P=n^{-1}I_n\). The exact square root therefore
prevents the variance from vanishing as the rank grows.
Lemma~\ref{lem:square-root-trace} simultaneously controls its
change under perturbation. The proof below replaces this exact
square root by a polynomial with an explicitly controlled error.

\begin{lemma}[From moments to bounded continuous functions]\label{lem:quadratic-moment-transfer}
Let \(X_j=X_j^*\) belong to unital \(\mathrm{C}^*\)-algebras
equipped with states \(\varphi_j\). Let \(q_B\) be a centered
Gaussian quadratic as in Lemma~\ref{lem:gaussian-quadratic}.
Suppose that
\[
 \varphi_j(X_j^d)\longrightarrow\mathbb E(q_B^d)
 \qquad(d=0,1,2,\ldots).
\]
Then, for every bounded continuous function \(f:\R\to\C\),
\[
 \varphi_j(f(X_j))\longrightarrow\mathbb E f(q_B).
\]
No uniform bound on \(\|X_j\|\) is needed.
\end{lemma}

\begin{proof}
Let \(\mu_j\) be the spectral probability law of \(X_j\) under
\(\varphi_j\). The second moments of \(\mu_j\) are bounded, so
the sequence is tight. Fix an integer \(d\geq0\), and choose an
even integer \(2r>d\). Moment convergence supplies a constant
\(C_r<\infty\) such that
\(\int |x|^{2r}\,d\mu_j(x)\leq C_r\) for all \(j\). Consequently
\begin{equation}\label{eq:spectral-moment-tail}
 \int_{|x|>K}|x|^d\,d\mu_j(x)
 \leq C_r K^{d-2r}
 \qquad(K>0).
\end{equation}

Consider any weakly convergent subsequence, with limit \(\mu\).
Lower semicontinuity gives
\(\int |x|^{2r}\,d\mu(x)\leq C_r\), so the same tail bound holds
for \(\mu\). Multiply \(x^d\) by a continuous cutoff which equals
one on \([-K,K]\) and zero outside \([-2K,2K]\). Weak convergence
applies to this bounded continuous function. Letting \(K\) tend
to infinity and using Equation~\eqref{eq:spectral-moment-tail} shows that
\[
 \int x^d\,d\mu(x)=\mathbb E(q_B^d).
\]
This holds for every \(d\). Lemma~\ref{lem:gaussian-quadratic}
therefore identifies \(\mu\) with the law of \(q_B\).
Tightness now shows that the whole sequence \(\mu_j\) converges
weakly to that law: every subsequence has a weakly convergent
further subsequence, and all such limits have just been
identified. This is precisely the claimed convergence for \(f\).
\end{proof}

\begin{proof}[Proof of Proposition~\ref{prop:uniform-unitaries}]
\textbf{Step 1. Choice of parameters.}
Fix \(H,\eps,m\) as in the statement. Choose \(t>0\) so that
\begin{equation}\label{eq:gaussian-t-choice}
 (1+t^2)^{-1/4}<\eps/4.
\end{equation}
Next choose \(\delta>0\) and \(0<\eta<1/2\) so that
\begin{equation}\label{eq:gaussian-error-choice}
 t\sqrt{2}\,(\sqrt{\delta}+2\eta)<\eps/4.
\end{equation}
Amenability gives a nonempty finite \(F\subset G\) satisfying
\[
 |hF\mathbin{\triangle}F|<\delta|F|
 \qquad(h\in H).
\]
Set
\[
 n=|F|,
 \qquad J=F\cup\bigcup_{h\in H}hF.
\]
Finally, choose a real polynomial
\begin{equation}\label{eq:square-root-polynomial}
 p(s)=\sum_{r=1}^{D}c_rs^r,
 \qquad
 \sup_{0\leq s\leq1}|p(s)-\sqrt{s}|<\eta/\sqrt n.
\end{equation}
Such a polynomial exists: first approximate \(\sqrt{s}\) on
\([0,1]\) by a real polynomial to within
\(\eta/(2\sqrt n)\), then subtract its value at zero.
Thus \(t,\delta,\eta,F,n,J\), and \(p\), including its degree
and coefficients, are all fixed before \(N\) tends to infinity.
None is chosen using a trace or a limiting covariance.

\textbf{Step 2. Uniform Gaussian comparison.}
Let \(R=(R_{g,g'})_{g,g'\in J}\) be any real positive semidefinite
matrix with diagonal one. Choose unit vectors \(v_g\) with Gram
matrix \(R\), spanning a finite-dimensional real Euclidean space
\(L\), and let \(W\) be standard Gaussian in \(L\). Write
\(\xi_g=\langle v_g,W\rangle\). For \(E=F\) or \(E=hF\), define
\[
 P_E=\frac1n\sum_{g\in E}v_gv_g^{\mathsf T},
 \qquad
 B_E=p(P_E),
 \qquad
 q_E=\langle W,B_EW\rangle-\Tr(B_E).
\]
The operator \(P_E\) is positive with trace one, so its spectrum
lies in \([0,1]\); its rank is at most \(n\).
Because \(p(0)=0\), the operator
\(p(P_E)-P_E^{1/2}\) vanishes on the kernel of \(P_E\).
The scalar approximation in Equation~\eqref{eq:square-root-polynomial}
therefore gives
\begin{equation}\label{eq:gaussian-rank-control}
 \|B_E-P_E^{1/2}\|_{\HS}^2
 \leq \operatorname{rank}(P_E)\,\eta^2/n
 \leq\eta^2,
 \qquad
 \|P_E^{1/2}\|_{\HS}^2=\Tr(P_E)=1.
\end{equation}
In particular,
\begin{equation}\label{eq:gaussian-nondegeneracy}
 1-\eta\leq\|B_E\|_{\HS}\leq1+\eta.
\end{equation}

The summands indexed by \(F\cap hF\) cancel when subtracting
\(P_F\) and \(P_{hF}\). Each remaining \(v_gv_g^{\mathsf T}\)
has trace norm one. Hence Lemma~\ref{lem:square-root-trace}
implies
\[
 \|P_F^{1/2}-P_{hF}^{1/2}\|_{\HS}^2
 \leq\|P_F-P_{hF}\|_1
 \leq\frac{|F\mathbin{\triangle}hF|}{n}<\delta.
\]
Combining this with Equation~\eqref{eq:gaussian-rank-control} yields
\[
 \|B_F-B_{hF}\|_{\HS}\leq\sqrt{\delta}+2\eta.
\]
Both forms \(q_F,q_{hF}\) are evaluated on the same \(W\), so
their difference is the centered quadratic associated with
\(B_{hF}-B_F\). Lemma~\ref{lem:gaussian-quadratic} gives
\begin{equation}\label{eq:gaussian-translation-bound}
 \|q_{hF}-q_F\|_{L^2}
 \leq\sqrt{2}\,(\sqrt{\delta}+2\eta).
\end{equation}
Also, Equation~\eqref{eq:gaussian-nondegeneracy} and \(\eta<1/2\)
give, for every integer \(\ell\geq1\),
\begin{equation}\label{eq:gaussian-fourier-bound}
 \bigl|\mathbb E\exp(\iu\ell t q_F)\bigr|
 \leq
 \bigl(1+4\ell^2t^2\|B_F\|_{\HS}^2\bigr)^{-1/4}
 \leq(1+t^2)^{-1/4}.
\end{equation}
All these bounds hold for every such \(R\), including singular
matrices. No invariance of \(R\) under translation of its indices
has been used.

\textbf{Step 3. Equivariant polynomial prescription.}
For \(r\geq1\), direct multiplication of the rank-one summands
gives
\[
 P_E^r
 =\frac1{n^r}
  \sum_{g_1,\ldots,g_r\in E}
  R_{g_1,g_2}\cdots R_{g_{r-1},g_r}
  v_{g_1}v_{g_r}^{\mathsf T}.
\]
Taking its quadratic form and its trace, respectively, shows that
\begin{equation}\label{eq:gaussian-covariance-polynomial}
 q_E=
 \sum_{r=1}^{D}\frac{c_r}{n^r}
 \sum_{g_1,\ldots,g_r\in E}
 R_{g_1,g_2}\cdots R_{g_{r-1},g_r}
 \bigl(\xi_{g_1}\xi_{g_r}-R_{g_r,g_1}\bigr).
\end{equation}
The product of covariance entries is interpreted as \(1\) when
\(r=1\).

Use the same prescription in \(C\) by setting
\begin{align}
 \widetilde Q_{N,E}
 &=
 \sum_{r=1}^{D}\frac{c_r}{n^r}
 \sum_{g_1,\ldots,g_r\in E}
 k_{N;g_1,g_2}\cdots k_{N;g_{r-1},g_r}
 \bigl(z_{N,g_1}z_{N,g_r}-k_{N;g_r,g_1}\bigr),
 \label{eq:empirical-quadratic-polynomial}\\
 Q_{N,E}
 &=\tfrac12\bigl(\widetilde Q_{N,E}
                 +\widetilde Q_{N,E}^{\,*}\bigr).
 \label{eq:empirical-quadratic-selfadjoint}
\end{align}
Multiplication in Equation~\eqref{eq:empirical-quadratic-polynomial} is
in the displayed order. The entries of the empirical covariance
array need satisfy no matrix positivity condition: this is a
finite polynomial in elements of \(C\), whereas square-root
functional calculus was applied only to \(P_E\) in the Gaussian
comparison.

The exact relabeling in Equation~\eqref{eq:empirical-equivariance} carries
each tuple in the sum for \(F\) to its left translate in the sum
for \(hF\). It preserves the order of
all factors, and taking self-adjoint parts commutes with
\(\bar\alpha_h\). Consequently
\begin{equation}\label{eq:quadratic-exact-equivariance}
 Q_{N,hF}=\bar\alpha_h(Q_{N,F})
 \qquad(h\in H,\ N\geq1).
\end{equation}

\textbf{Step 4. Transfer of moments and scalar laws.}
Consider now any sequence \(N_j\to\infty\) and extreme traces
\(\varphi_j\in\T(C)\) such that
\[
 \varphi_j(k_{N_j;g,g'})\longrightarrow R_{g,g'}
 \qquad(g,g'\in J).
\]
Proposition~\ref{prop:joint-clt} gives convergence of every
ordered mixed moment of the empirical variables. Any moment of
the polynomials in
Equations~\eqref{eq:empirical-quadratic-polynomial}--\eqref{eq:empirical-quadratic-selfadjoint}
is a finite linear combination of such moments. Since the
limiting variables commute and
Equation~\eqref{eq:gaussian-covariance-polynomial} is real-valued, it
follows that the \(Q_{N_j,E}\) converge jointly in moments under
\(\varphi_j\) to the \(q_E\). In particular,
\begin{equation}\label{eq:quadratic-difference-limit}
 \|Q_{N_j,hF}-Q_{N_j,F}\|_{2,\varphi_j}
 \longrightarrow \|q_{hF}-q_F\|_{L^2}.
\end{equation}
Lemma~\ref{lem:quadratic-moment-transfer} also gives
\begin{equation}\label{eq:quadratic-exponential-limit}
 \varphi_j\bigl(\exp(\iu s Q_{N_j,F})\bigr)
 \longrightarrow\mathbb E\exp(\iu s q_F)
 \qquad(s\in\R).
\end{equation}
This use of scalar spectral laws is pertinent because the
operator norms of the \(Q_{N,E}\) may grow with \(N\).
Indeed, the elementary estimates
\(\|k_{N;g,g'}\|\leq M^2\) and
\(\|z_{N,g}\|\leq\sqrt N\,M\) yield only
\[
 \|Q_{N,E}\|\leq(N+1)\sum_{r=1}^{D}|c_r|M^{2r}.
\]
Convergence of the scalar moments, together with the tail
estimate Equation~\eqref{eq:spectral-moment-tail}, is sufficient for
Equation~\eqref{eq:quadratic-exponential-limit}.

\textbf{Step 5. Uniformity over all traces.}
Set \(V_N=\exp(\iu tQ_{N,F})\). For any self-adjoint \(X,Y\)
in a unital \(\mathrm{C}^*\)-algebra and any tracial state
\(\varphi\), differentiation gives
\[
 \frac{d}{ds}
 \bigl(e^{\iu(1-s)tX}e^{\iu stY}\bigr)
 =e^{\iu(1-s)tX}\,\iu t(Y-X)\,e^{\iu stY}.
\]
The tracial \(2\)-seminorm is unchanged by left and right
unitary multiplication. Integration and the triangle inequality
therefore give
\begin{equation}\label{eq:tracial-duhamel}
 \|e^{\iu tX}-e^{\iu tY}\|_{2,\varphi}
 \leq t\|X-Y\|_{2,\varphi}.
\end{equation}
In view of Equations~\eqref{eq:quadratic-exact-equivariance},
\eqref{eq:quadratic-difference-limit}, and
\eqref{eq:gaussian-translation-bound}, every sequence of traces
and indices with convergent covariance as above satisfies
\begin{align}
 \limsup_{j\to\infty}
 \|\bar\alpha_h(V_{N_j})-V_{N_j}\|_{2,\varphi_j}
 &\leq t\sqrt{2}\,(\sqrt{\delta}+2\eta)<\eps/4,
 \label{eq:sequential-invariance-estimate}\\
 \lim_{j\to\infty}|\varphi_j(V_{N_j}^{\ell})|
 &=
 \bigl|\mathbb E\exp(\iu\ell t q_F)\bigr|
 \leq(1+t^2)^{-1/4}<\eps/4.
 \label{eq:sequential-fourier-estimate}
\end{align}
Here \(h\in H\) and \(1\leq\ell\leq m\), and the strict
inequalities are Equations~\eqref{eq:gaussian-error-choice} and
\eqref{eq:gaussian-t-choice}.

It follows that, for all sufficiently large \(N\), every extreme
trace \(\varphi\in\T(C)\) satisfies
\begin{equation}\label{eq:extreme-uniform-estimates}
 \|\bar\alpha_h(V_N)-V_N\|_{2,\varphi}<\eps/2
 \quad(h\in H),\qquad
 |\varphi(V_N^\ell)|<\eps/2
 \quad(1\leq\ell\leq m).
\end{equation}
Indeed, otherwise choose \(N_j\to\infty\) and an extreme trace
\(\varphi_j\) for which one of these finitely many tests fails.
Passing to a subsequence fixes the failed test. The matrices
\[
 \bigl(\varphi_j(k_{N_j;g,g'})\bigr)_{g,g'\in J}
\]
belong to the compact set of real positive semidefinite matrices
with diagonal one. A further subsequence has a limiting
covariance \(R\), and
Equation~\eqref{eq:sequential-invariance-estimate} or
Equation~\eqref{eq:sequential-fourier-estimate} contradicts the failure
at tolerance \(\eps/2\).
Only the finite covariance matrices were extracted; no
compactness of the extreme trace boundary is needed. Nor is a
convergence rate uniform in \(R\) required: each putative
sequence of failures supplies one limiting \(R\), to which the
preceding moment argument applies. Moment determinacy is used
after that \(R\) has been fixed, so no common exponential-moment
radius over all covariance matrices is required.

Fix an \(N\) for which Equation~\eqref{eq:extreme-uniform-estimates}
holds. For each \(h\in H\), the squared seminorm
\[
 \varphi\longmapsto
 \varphi\bigl((\bar\alpha_h(V_N)-V_N)^*
              (\bar\alpha_h(V_N)-V_N)\bigr)
\]
is weak-star continuous and affine on \(\T(C)\). For each
\(1\leq\ell\leq m\), the map
\(\varphi\mapsto|\varphi(V_N^\ell)|\) is weak-star continuous
and convex. The sublevel sets with bounds \(\eps^2/4\) and
\(\eps/2\), respectively, are closed and convex and contain
all extreme traces. By the Krein--Milman Theorem they contain
\(\T(C)\). Thus the suprema in
Equations~\eqref{eq:uniform-unitary-invariance} and
\eqref{eq:uniform-unitary-moments} are at most \(\eps/2\),
and in particular strictly less than \(\eps\).
Taking \(Q=Q_{N,F}\) and \(V=V_N\) completes the proof.
\end{proof}
